% Options for packages loaded elsewhere
% Options for packages loaded elsewhere
\PassOptionsToPackage{unicode}{hyperref}
\PassOptionsToPackage{hyphens}{url}
\PassOptionsToPackage{dvipsnames,svgnames,x11names}{xcolor}
%
\documentclass[
  german,
  10pt,
  twoside]{article}
\usepackage{xcolor}
\usepackage[top=2.5cm,bottom=2.5cm,left=2cm,right=2cm]{geometry}
\usepackage{amsmath,amssymb}
\setcounter{secnumdepth}{5}
\usepackage{iftex}
\ifPDFTeX
  \usepackage[T1]{fontenc}
  \usepackage[utf8]{inputenc}
  \usepackage{textcomp} % provide euro and other symbols
\else % if luatex or xetex
  \usepackage{unicode-math} % this also loads fontspec
  \defaultfontfeatures{Scale=MatchLowercase}
  \defaultfontfeatures[\rmfamily]{Ligatures=TeX,Scale=1}
\fi
\usepackage{lmodern}
\ifPDFTeX\else
  % xetex/luatex font selection
\fi
% Use upquote if available, for straight quotes in verbatim environments
\IfFileExists{upquote.sty}{\usepackage{upquote}}{}
\IfFileExists{microtype.sty}{% use microtype if available
  \usepackage[]{microtype}
  \UseMicrotypeSet[protrusion]{basicmath} % disable protrusion for tt fonts
}{}
\usepackage{setspace}
\makeatletter
\@ifundefined{KOMAClassName}{% if non-KOMA class
  \IfFileExists{parskip.sty}{%
    \usepackage{parskip}
  }{% else
    \setlength{\parindent}{0pt}
    \setlength{\parskip}{6pt plus 2pt minus 1pt}}
}{% if KOMA class
  \KOMAoptions{parskip=half}}
\makeatother
% Make \paragraph and \subparagraph free-standing
\makeatletter
\ifx\paragraph\undefined\else
  \let\oldparagraph\paragraph
  \renewcommand{\paragraph}{
    \@ifstar
      \xxxParagraphStar
      \xxxParagraphNoStar
  }
  \newcommand{\xxxParagraphStar}[1]{\oldparagraph*{#1}\mbox{}}
  \newcommand{\xxxParagraphNoStar}[1]{\oldparagraph{#1}\mbox{}}
\fi
\ifx\subparagraph\undefined\else
  \let\oldsubparagraph\subparagraph
  \renewcommand{\subparagraph}{
    \@ifstar
      \xxxSubParagraphStar
      \xxxSubParagraphNoStar
  }
  \newcommand{\xxxSubParagraphStar}[1]{\oldsubparagraph*{#1}\mbox{}}
  \newcommand{\xxxSubParagraphNoStar}[1]{\oldsubparagraph{#1}\mbox{}}
\fi
\makeatother


\usepackage{graphicx}
\makeatletter
\newsavebox\pandoc@box
\newcommand*\pandocbounded[1]{% scales image to fit in text height/width
  \sbox\pandoc@box{#1}%
  \Gscale@div\@tempa{\textheight}{\dimexpr\ht\pandoc@box+\dp\pandoc@box\relax}%
  \Gscale@div\@tempb{\linewidth}{\wd\pandoc@box}%
  \ifdim\@tempb\p@<\@tempa\p@\let\@tempa\@tempb\fi% select the smaller of both
  \ifdim\@tempa\p@<\p@\scalebox{\@tempa}{\usebox\pandoc@box}%
  \else\usebox{\pandoc@box}%
  \fi%
}
% Set default figure placement to htbp
\def\fps@figure{htbp}
\makeatother



\ifLuaTeX
\usepackage[bidi=basic]{babel}
\else
\usepackage[bidi=default]{babel}
\fi
% get rid of language-specific shorthands (see #6817):
\let\LanguageShortHands\languageshorthands
\def\languageshorthands#1{}
\ifLuaTeX
  \usepackage[german]{selnolig} % disable illegal ligatures
\fi


\setlength{\emergencystretch}{3em} % prevent overfull lines

\providecommand{\tightlist}{%
  \setlength{\itemsep}{0pt}\setlength{\parskip}{0pt}}



 
\usepackage[style=authoryear,backend=biber,style=authoryear-comp,uniquename=false,uniquelist=false]{biblatex}
\addbibresource{../references.bib}


\usepackage{booktabs}
\usepackage{longtable}
\usepackage{array}
\usepackage{multirow}
\usepackage{wrapfig}
\usepackage{float}
\usepackage{colortbl}
\usepackage{pdflscape}
\usepackage{tabu}
\usepackage{threeparttable}
\usepackage{threeparttablex}
\usepackage[normalem]{ulem}
\usepackage{makecell}
\usepackage{xcolor}
\usepackage{csquotes}
\usepackage{booktabs}
\usepackage{array}
\usepackage{multirow}
\usepackage{wrapfig}
\usepackage{float}
\usepackage{colortbl}
\usepackage{pdflscape}
\usepackage{threeparttable}
\usepackage{threeparttablex}
\usepackage[normalem]{ulem}
\usepackage{makecell}
\usepackage{xcolor}
\usepackage{multicol}
\usepackage{fancyhdr}
\usepackage{titlesec}
\usepackage{etoolbox}
\usepackage{tabularx}
\usepackage{afterpage}
\usepackage{tikz}
\usepackage{longtable}



% Suppress default maketitle
\renewcommand{\maketitle}{}

% Define WISTA colors (Havelock Blue as main)
\definecolor{wistaBlue}{RGB}{65,125,180}
\definecolor{wistaLightBlue}{RGB}{120,170,215}
\definecolor{wistaDarkGray}{RGB}{64,64,64}
\definecolor{wistaLightGray}{RGB}{128,128,128}

% Section formatting (##) - number 20% bigger, lines closer, text closer to bottom line
% e.g., "2" with lines around "Einleitung"
\titleformat{\section}[display]
  {\normalfont\large\bfseries}
  {\color{wistaBlue}\Large\bfseries\thesection\\\noindent\rule{\columnwidth}{1.5pt}}
  {0pt}
  {\vspace{0pt}\color{wistaDarkGray}\large\bfseries}
  [\vspace{0pt}\noindent{\color{wistaBlue}\rule{\columnwidth}{1.5pt}}\vspace{0.05cm}]
\titlespacing*{\section}{0pt}{1.5ex plus 0.3ex minus .1ex}{0.2ex plus .1ex}

% Subsection formatting (###) - grey text with blue line under
% e.g., "2.1 Handlungsbedarf"
\renewcommand{\thesubsection}{\thesection.\arabic{subsection}}
\titleformat{\subsection}
  {\color{wistaDarkGray}\normalsize\bfseries}
  {\color{wistaDarkGray}\normalsize\bfseries\thesubsection}
  {0.5em}
  {\color{wistaDarkGray}\normalsize\bfseries}
  [\vspace{0pt}\noindent{\color{wistaBlue}\rule{\columnwidth}{0.8pt}}\vspace{0.05cm}]
\titlespacing*{\subsection}{0pt}{1ex plus 0.2ex minus .1ex}{0.2ex plus .1ex}
  
% Subsubsection formatting (####) - grey text, no underline
% e.g., "2.1.1 Steigende Datenmenge"
\renewcommand{\thesubsubsection}{\thesubsection.\arabic{subsubsection}}
\titleformat{\subsubsection}
  {\color{wistaDarkGray}\normalsize\bfseries}
  {\color{wistaDarkGray}\normalsize\bfseries\thesubsubsection}
  {0.5em}
  {\color{wistaDarkGray}\normalsize\bfseries}

% Paragraph formatting (4th level) - numbered as subsection.subsubsection.paragraph
\renewcommand{\theparagraph}{\thesubsubsection.\arabic{paragraph}}
\titleformat{\paragraph}[runin]
  {\color{wistaDarkGray}\normalsize\bfseries}
  {\color{wistaBlue}\normalsize\bfseries\theparagraph}
  {0.5em}
  {\color{wistaDarkGray}\normalsize\bfseries}
  [.]
\titlespacing*{\paragraph}{0pt}{1ex plus 0.5ex minus 0.2ex}{0.5em}

% Header and footer setup with blue line under header
\pagestyle{fancy}
\fancyhf{}
\fancyhead[L]{\small\textit{Oliver Hauke, Annette Kristiansen}}
\fancyhead[R]{\small\textit{Effiziente Analyse großer Forschungsdaten}}
\fancyfoot[L]{\small\thepage}
\fancyfoot[R]{\small Statistisches Bundesamt | WISTA | {\color{wistaBlue}6} | 2025}
\renewcommand{\headrulewidth}{0pt}
\renewcommand{\footrulewidth}{0pt}
% Add blue line under header
\renewcommand{\headrule}{\color{wistaBlue}\hrule width\headwidth height 1pt \vskip 2pt}
\renewcommand{\headrulewidth}{1pt}

% Make all figures span both columns (full width)
\let\origfigure\figure
\let\endorigfigure\endfigure
\renewenvironment{figure}[1][htbp]{%
  \origfigure*[#1]%
}{%
  \endorigfigure*%
}

% Allow manual control of table* environment for full-width tables
% Tables by default stay in single column, but can use table* manually



% Make bibliography full-width (one column) - for biblatex
\defbibheading{bibliography}[\bibname]{%
  \onecolumn
  \section*{#1}%
  \markboth{#1}{#1}%
}



% Footnotes stay within column
\usepackage[stable]{footmisc}

% Custom footnote formatting: "|^1" in dark WISTA gray
\renewcommand{\thefootnote}{{\color{wistaDarkGray}|$^{\arabic{footnote}}$}}

% Pifont for dingbat symbols
\usepackage{pifont}

% Better table handling for two-column mode
\usepackage{array,tabularx,calc}

% Custom WISTA table format
\usepackage{caption}
\usepackage{xcolor}
\usepackage{colortbl}
\usepackage{array}

% Define custom caption format for WISTA tables
\DeclareCaptionFormat{wista}{%
  {\color{wistaLightBlue}\bfseries #1#2}\par\vspace{0.2em}%
  {\color{wistaDarkGray}\normalfont #3}\par\vspace{0.3em}%
}

% Set up table captions with custom WISTA format
\captionsetup[table]{
  position=top,
  format=wista,
  justification=justified,
  singlelinecheck=false,
  labelsep=none,
  skip=0.3em
}

% Better table spacing and white inner lines
\setlength{\tabcolsep}{6pt}
\renewcommand{\arraystretch}{1.15}

% Set white lines for tables by default, blue for first column
\AtBeginDocument{
  \arrayrulecolor{white}
  \setlength{\arrayrulewidth}{0.3pt}
}

% Define commands for colored table borders
\newcommand{\bluevrule}{\color{wistaBlue}\vrule}
\newcommand{\whitevrule}{\color{white}\vrule}

% Define WISTA table colors
\definecolor{wistaTableBg}{RGB}{240,248,255}     % Very light blue background
\definecolor{wistaHeaderBg}{RGB}{220,230,240}    % Light blue header background
\definecolor{wistaFirstCol}{RGB}{65,125,180}     % Blue for first column lines

% Custom WISTA table command
\newcommand{\wistaTableStart}{%
  \rowcolors{2}{wistaTableBg}{wistaTableBg}%
  \arrayrulecolor{wistaFirstCol}%
}

% Custom commands for table styling
\newcommand{\wistaHeaderRow}[1]{%
  \rowcolor{wistaHeaderBg}%
  \arrayrulecolor{wistaDarkGray}%
  #1 \\
  \arrayrulecolor{white}%
}

\newcommand{\wistaFirstColRule}{\arrayrulecolor{wistaFirstCol}}
\newcommand{\wistaWhiteRule}{\arrayrulecolor{white}}

% Custom bullet points with > character - 35% width, 95% height, black extra bold, half line spacing
\usepackage{enumitem}
\setlist[itemize]{label={\scalebox{0.35}[0.95]{\fontseries{bx}\selectfont >}}, leftmargin=1.5em, itemsep=5pt, topsep=3pt, parsep=0pt}

% Redefine ref to include pifont 247 (arrow) for figure/table/section references
\let\oldref\ref
\renewcommand{\ref}[1]{\oldref{#1}\raisebox{-0.2ex}{\scriptsize\color{wistaBlue}\ding{247}}}
\makeatletter
\@ifpackageloaded{caption}{}{\usepackage{caption}}
\AtBeginDocument{%
\ifdefined\contentsname
  \renewcommand*\contentsname{Inhaltsverzeichnis}
\else
  \newcommand\contentsname{Inhaltsverzeichnis}
\fi
\ifdefined\listfigurename
  \renewcommand*\listfigurename{Abbildungsverzeichnis}
\else
  \newcommand\listfigurename{Abbildungsverzeichnis}
\fi
\ifdefined\listtablename
  \renewcommand*\listtablename{Tabellenverzeichnis}
\else
  \newcommand\listtablename{Tabellenverzeichnis}
\fi
\ifdefined\figurename
  \renewcommand*\figurename{Abbildung}
\else
  \newcommand\figurename{Abbildung}
\fi
\ifdefined\tablename
  \renewcommand*\tablename{Tabelle}
\else
  \newcommand\tablename{Tabelle}
\fi
}
\@ifpackageloaded{float}{}{\usepackage{float}}
\floatstyle{ruled}
\@ifundefined{c@chapter}{\newfloat{codelisting}{h}{lop}}{\newfloat{codelisting}{h}{lop}[chapter]}
\floatname{codelisting}{Listing}
\newcommand*\listoflistings{\listof{codelisting}{Listingverzeichnis}}
\makeatother
\makeatletter
\makeatother
\makeatletter
\@ifpackageloaded{caption}{}{\usepackage{caption}}
\@ifpackageloaded{subcaption}{}{\usepackage{subcaption}}
\makeatother
\usepackage{bookmark}
\IfFileExists{xurl.sty}{\usepackage{xurl}}{} % add URL line breaks if available
\urlstyle{same}
\hypersetup{
  pdftitle={EFFIZIENTE AUSWERTUNG DES BUSINESS-TAX-PANELS MITHILFE VON R},
  pdfauthor={Oliver Hauke; Annette Kristiansen},
  pdflang={de},
  colorlinks=true,
  linkcolor={blue},
  filecolor={Maroon},
  citecolor={Blue},
  urlcolor={Blue},
  pdfcreator={LaTeX via pandoc}}


\title{EFFIZIENTE AUSWERTUNG DES BUSINESS-TAX-PANELS MITHILFE VON R}
\author{Oliver Hauke \and Annette Kristiansen}
\date{2025-12-04}
\begin{document}
\maketitle

% Custom WISTA-style title page
\begin{titlepage}
\thispagestyle{empty}
\pagestyle{empty}
\onecolumn

% Horizontal blue line at top (header line)
\noindent{\color{wistaBlue}\rule{\textwidth}{1pt}}

\vspace{-0.5cm}

% Vertical blue line on the left (6.3cm from left edge, starts below header line, stops above footer)
\begin{tikzpicture}[remember picture, overlay]
  \draw[wistaBlue, line width=2pt] 
    ([xshift=6.3cm, yshift=-4cm]current page.north west) -- 
    ([xshift=6.3cm, yshift=3.5cm]current page.south west);
\end{tikzpicture}

\vspace*{0.5cm}

% Left side - short author bios (right-aligned to the vertical line)
\begin{minipage}[t]{3.9cm}
\selectlanguage{ngerman}%
\raggedleft
\hyphenpenalty=0
\tolerance=1000
\fontsize{7}{8.5}\selectfont
\vspace{0pt}%
{\fontsize{8}{9.5}\selectfont\color{wistaBlue}\textbf{Oliver Hauke}}\\[0.15cm]
ist Mathematiker und IT-Referent im Kompetenzzentrum „Analyse und Auswertung" des Statistischen Bundesamtes. Er verantwortet die Weiterentwicklung der technischen Infrastruktur des Forschungsdatenzentrums und berät das Netzwerk empirische Steuerforschung in der effizienten Nutzung der Systeme.\\[0.5cm]

{\fontsize{8}{9.5}\selectfont\color{wistaBlue}\textbf{Annette Kristiansen}}\\[0.15cm]
ist Ökonomin und Referentin im Referat „Unternehmenssteuern" des Statistischen Bundesamtes. Sie beschäftigt sich mit der Zusammenführung der Unternehmenssteuerstatistiken und betreut das Netzwerk empirische Steuerforschung. Für ihre externe Promotion an der Universität Bremen untersucht sie die technologische Vielfalt multinationaler Unternehmen.

\end{minipage}%
\hspace{0.4cm}% Gap to vertical line (left side)
\hspace{0.4cm}% Gap from vertical line (right side)
% Right side - title, authors, keywords, abstracts
\begin{minipage}[t]{11.8cm}
\raggedright
\vspace{1.2cm}%

% Title (uppercase via command, not bold)
{\LARGE\color{wistaDarkGray}\begin{spacing}{1.3}\MakeUppercase{Effiziente Auswertung des Business-Tax-Panels mithilfe von R}\end{spacing}}

\vspace{-0.2cm}

% Blue horizontal line (underline for title)
{\color{wistaBlue}\rule{\linewidth}{0.5pt}}

\vspace{0.5cm}

% Authors inline, dark gray
{\large\color{wistaDarkGray} Oliver Hauke, Annette Kristiansen}

\vspace{0.05cm}

% Blue line (underline for authors)
{\color{wistaBlue}\rule{\linewidth}{0.5pt}}

\vspace{0.5cm}

% Keywords (pifont 247 arrow then bold label and blue text)
\hangindent=1.2em\hangafter=1
{\color{wistaBlue}\ding{247}\ {\bfseries Schlüsselwörter:} Effiziente Analysen, R-Programmierung, Forschungsdaten, Unternehmenssteuerstatistik, Business-Tax-Panel, Netzwerk empirische Steuerforschung}

\vspace{0.3cm}

% Zusammenfassung (uppercase via command)
{\color{wistaDarkGray}\bfseries\MakeUppercase{Zusammenfassung}}\\[0.15cm]
\small
Das Forschungsdatenzentrum des Statistischen Bundesamtes baut seine technische Infrastruktur gezielt aus, um der Wissenschaft erweiterte Analysemöglichkeiten amtlicher Daten bereitzustellen. Seit November 2025 steht Forschenden exklusiver Zugang zu diesen erweiterten Systemressourcen in R und Stata zur Verfügung. R-basierte Tools -- insbesondere Apache Arrow, Parquet und DuckDB -- ermöglichen es, gängige Analysen großer Forschungsdatensätze in Echtzeit auszuführen. Am Beispiel des Business-Tax-Panel (2013 bis 2019) werden die erzielbaren Effizienzgewinne bei der Auswertung umfangreicher steuerstatistischer Daten im Rahmen des Netzwerks empirische Steuerforschung demonstriert.

\vspace{0.5cm}

% English keywords and abstract (italicized)
\hangindent=1.2em\hangafter=1
{\itshape\color{wistaBlue}\ding{247}\ {\bfseries Keywords:} Efficient analysis -- R programming -- research data -- corporate tax statistics -- Business Tax Panel -- empirical tax research network}

\vspace{0.3cm}

{\itshape\color{wistaDarkGray}\bfseries\MakeUppercase{Abstract}}\\[0.15cm]
\small\itshape
The Research Data Centre at the Federal Statistical Office is strategically expanding its technical infrastructure to provide researchers with enhanced analytical capabilities for official data. Since November 2025, researchers have had exclusive access to these extended system resources in R and Stata. R-based tools -- particularly Apache Arrow, Parquet, and DuckDB -- enable real-time analysis of large research datasets. Using the Business Tax Panel (2013--2019), this paper demonstrates the achievable efficiency gains in analysing large-scale tax statistics within the empirical tax research network.

\end{minipage}

\vfill

\vspace{0.5cm}
% Footer matching document style
\noindent\makebox[\textwidth]{%
  \small 1%
  \hfill%
  \small Statistisches Bundesamt | WISTA | {\color{wistaBlue}6} | 2025%
}

\end{titlepage}

% Stay in one-column mode for TOC (will be switched later)
\pagestyle{fancy}

\renewcommand*\contentsname{Inhaltsverzeichnis}
{
\hypersetup{linkcolor=}
\setcounter{tocdepth}{3}
\tableofcontents
}

\setstretch{1}
\section{Datenverarbeitungsmethoden}\label{sec-methoden}

\subsection{Ausgangslage}\label{ausgangslage}

Wir betrachten in dieser Arbeit die Performanz sog.
Datenverarbeitungsmethoden, d.h. Softwarelösungen, die Daten
einlesen,modifzierern, aggregieren, auswerten oder analysieren. Das
Schreiben von Daten wird in unserer Untersuchung ausgeklammert, da es
sich bei den Veröffentlichungen des Statistischen Bundesamtes oder den
Outputs der Wissenschaft im FDZ um kleine Tabellen handelt, die keine
große Performanz erfordern. Dennoch bieten die hier vorstellten Lösungen
auch hohe Performanz im Schreiben von Daten.\footnote{\begin{itemize}
  \tightlist
  \item
    Hauptanalysewerkzeuge sind bisher:

    \begin{itemize}
    \tightlist
    \item
      im StBA SAS 9.4 M8, EG
    \item
      im FDZ zsl. Stata
    \end{itemize}
  \item
    bisher: V.15 MP-12 (KDFV) und Stata 19 SE lokal (GWAP).
  \item
    neu im OSS in KDFV/GWAP V. 19 gleichermaßen in großen
    Systemkapazitäten.
  \end{itemize}}

Daneben wächst der Einsatz von R stetig im StBA und FDZ stetig.

Empfehlung des FDSZ der deutschen Bundesbank
\textcite{gomolka2021largedata} \ldots.

\subsection{Anforderungen an
Datenverarbeitungsmethoden}\label{anforderungen-an-datenverarbeitungsmethoden}

\begin{itemize}
\tightlist
\item
  Performanz und Nutzererfahrung auf dem aktuellen Stand der Technik
\item
  leicht einzurichten auch in den abgeschotteten System der amtl.
  Statistik (tlw. langwierig und schwierig), d.h. wenig Abhängigkeiten
  und unabhängig von Internet bzw. Cloud betreibbar.
\item
  langfristig stabil und gut unterstützt
\item
  leicht nutzbar, gut dokumentiert, gutes Schulungsangebot
\end{itemize}

\subsection{Geeignete Datenverarbeitungsmethoden in
R}\label{geeignete-datenverarbeitungsmethoden-in-r}

\begin{itemize}
\tightlist
\item
  R ist die native Programmiersprache von Statistiker/-innen und
  Daten-affinen Programmierer/-innen. Frei verfügbar und leicht
  einzurichten. Somit weit verbreitet, insb. an Universitäten.
\item
  R ist Open Source und wird ständig weiterentwickelt (aktuell in
  Version 4.5 verwendet, jährliche Major Releases).
\item
  Eine weltweit aktive Community löst Probleme zeitnah und ergänzt
  stetig Funktionalitäten in Form von sog. \emph{R-Packages}.
\item
  Bietet nativ bereits viele Datenverarbeitungsmethoden, die jedoch
  durch diese R-Packages ganz neue Level an Nutzerfreundlichkeit,
  Funktionsumfang und Performanz erreichen.
\end{itemize}

\subsubsection{Klassische Methoden}\label{klassische-methoden}

\textbf{Base R.}

\textbf{Vorteile:} Einfach, weit verbreitet, keine zusätzlichen
Abhängigkeiten \textbf{Nachteile:} Langsam bei großen Daten, hoher
Speicherverbrauch, alle Daten müssen in RAM passen

\textbf{Empfehlung:} Nur für kleine Datensätze (\textless1 GB) oder wenn
Kompatibilität wichtiger als Performance ist

\begin{itemize}
\tightlist
\item
  base-r Grundlage/einstieg, weder sonderlich amschaulich nocb
  sonderlich performant -\textgreater{} Packages überlassen
\end{itemize}

Zweckdienlich für Einsteiger oder gelegenheitsnutzer oder falls gar
keine Abhängigkeiten in Form von Pkg erlaubt sind. Ansonsten\ldots{}

\textbf{data.table.}

\textbf{Vorteile:} Sehr schnell für In-Memory-Operationen, ausgereift,
kompakte Syntax \textbf{Nachteile:} Alle Daten müssen in RAM passen,
steile Lernkurve, keine Out-of-Core-Unterstützung

\textbf{Empfehlung:} Gute Alternative wenn Datensatz in RAM passt und
keine Parquet-Integration benötigt wird

-\textgreater{} \{data.table\} in den meisten Fällen die performanteste
klassische DVM. An die Syntax an base-r angesiedelt und nur hier
verwendet, spezielles Wissen notwendig.

\begin{itemize}
\tightlist
\item
  \{tidyverse\} Ökosystem von Packages zur Datenverarbeitung. bekannt
  für besonders hohe Nutzerfreundlichkeit und einer eigenen
  Programmiersyntax sog. \{tidy\}-Syntax, die den Anspruch hat
  menschlicher Sprache abzubilden:
\end{itemize}

{[} Tabelle Beispiel Dplyr {]}

\textbf{dplyr + dbplyr.}

\textbf{Vorteile:} Vertraute dplyr-Syntax, gut mit Datenbanken
integriert \textbf{Nachteile:} Langsamer als Arrow/DuckDB, benötigt
Backend-Datenbank, SQL-Translation manchmal limitiert

\textbf{Empfehlung:} Wenn bereits PostgreSQL- oder ähnliche
Datenbank-Infrastruktur vorhanden ist

So beliebt, dass die Community sie vielen anderen Packages ergänzt hat.
Bspw. für \{data.table\} die Packages \{dtplyr\} oder \{tidytable\}. Wie
im Falle von \{dtplyr\} kann es dabei aber zu Einschränkung der
Performamz kommen.\footnote{\ldots{}}

\subsubsection{Ausgewählte Best
Practices}\label{ausgewuxe4hlte-best-practices}

in Abschnitt Kapitel~\ref{sec-methoden} beschriebenen Methoden
(Variablenselektion, Datenaufteilung):

\subsubsection{Hoch-performante
R-Packages}\label{hoch-performante-r-packages}

Definition einfach: können nativ Daten über RAM verarbeiten.

von Desktopanwendung unterm Tisch zu Performamten Analyseserver
paradigmenwechsel notwendig: in RAM vs nicht im RAM!

Die auf den Servern des StBA und FDZ verfügbaren Packageversionen lassen
sich aus organisatorischen Gründen nicht flexibel anpassen (kein
Internet, kuration intern, s. A. Inf) können Tabelle X entnommen werden.

\begin{table*}

\caption{\label{tbl-pkg-packages}Versionen betrachteter R-Pakete}

\centering{

\centering\begingroup\fontsize{10}{12}\selectfont


\begin{tabular}{|>{}l|||>{}l|}
\hline
\cellcolor[HTML]{dce6f0}{\textcolor[HTML]{404040}{\textbf{Paket}}} & \cellcolor[HTML]{dce6f0}{\textcolor[HTML]{404040}{\textbf{Version}}}\\
\hline
\cellcolor{white}{\textcolor{black}{\{data.table\}}} & \cellcolor[HTML]{f0f8ff}{1.17.8}\\
\cellcolor{white}{\textcolor{black}{\{tidyverse\}}} & \cellcolor[HTML]{f0f8ff}{tba}\\
\cellcolor{white}{\textcolor{black}{\{vroom\}}} & \cellcolor[HTML]{f0f8ff}{tba}\\
\cellcolor{white}{\textcolor{black}{\{arrow\}}} & \cellcolor[HTML]{f0f8ff}{20.0.0.2}\\
\cellcolor{white}{\textcolor{black}{\{duckdb\}}} & \cellcolor[HTML]{f0f8ff}{1.3.2}\\
\cellcolor{white}{\textcolor{black}{\{polars\}}} & \cellcolor[HTML]{f0f8ff}{1.0.1}\\
\hline
\end{tabular}
\endgroup{}

}

\end{table*}%

Wie in Abschnitt \textbf{?@sec-infrastruktur} beschrieben war es uns aus
organisatorischen Gründen nicht möglich, die neuste Version aller Pakete
zu verwenden. Die leichter nutzbaren Variationen \{duckplyr\} bzw.
\{tidytable\} zu nutzen. Die Package wurden aber wenige Tage nach dem
Experiment eingerichtet. Die Syntaxtranslation kann dabei die Performanz
verschlechtern, sodass bei \{duckdb\} bzw. \{data.table\} zwischen
Performanz und Code-Verständlichkeit abgewägt werden muss. \{arrow\} ist
schon länger als \{polars\} im Statistischen Bundesamt verfügbar, sodass
wir darin bereits mehr Erfahrungen einsetzen konnten.

Die Analyse großer Datensätze wie des in der FDZ-Umgebung stellt
besondere Anforderungen an die verwendeten Werkzeuge. Traditionelle
Ansätze (CSV-Dateien, Base R) stoßen bei dieser Datengröße an ihre
Grenzen. Moderne spaltenorientierte Technologien bieten hier deutliche
Vorteile.

beachte/erwähne:
https://www.r-bloggers.com/2024/04/the-truth-about-tidy-wrappers-3/


\printbibliography



\end{document}
